\section{Introduction}
\label{sec:intro}

The Border Gateway Protocol (BGP) is the de facto inter-domain routing protocol of the Internet: it propagates route announcements across more than seventy-three thousand autonomous systems (ASes) and is responsible for telling each network which prefixes are reachable and through which path~\cite{rfc4271}. What makes BGP practical, however, is also what makes it fragile. BGP was designed without any built-in authentication or authorization mechanism, and its trust model assumes that every AS truthfully announces only the prefixes it owns or is entitled to transit. As a result, any AS can announce any prefix, and the global routing system must honour the announcement until the operator community manually intervenes. This absence of endogenous security has made BGP incidents frequent and costly: in August~2022 a cryptocurrency exchange lost \$2.3M in user funds when a BGP route leak diverted withdrawal traffic to a malicious AS for 90~minutes~\cite{krenc-2023-bgp-hijack}. Principled defences do exist, including cryptographic extensions (BGPsec, RPKI with ROV~\cite{cooper-2013-rpki}) and the MANRS operator practices~\cite{manrs}, but their deployment is incremental and the coordination overhead limits adoption at scale. Consequently, BGP security problems persist, and anomaly attribution remains the operator's responsibility for the foreseeable future.

In practice, this responsibility falls on operators who are already saturated. To detect routing anomalies, they rely on BGP monitoring platforms that ingest public telemetry and raise alerts whenever an announcement deviates from an expected baseline. These platforms are effective at surfacing something, but they surface far too much: a single operator may receive thousands of alerts per day, of which fewer than one percent are ultimately confirmed as actionable. The triage that stands between this stream and a decision still rests on static thresholds and heuristic rules, which cannot tell a routing anomaly that matters from one that does not. The consequence is an alert storm: genuinely high-impact incidents are buried in noise, and the operator's attention—the scarcest resource in the loop—is spent on events that never affect the business.

The root cause of this failure is not a lack of detectors, but a \textbf{disconnect between the routing signal and the traffic signal}. Existing work typically examines these signals in isolation: routing monitors characterize announcements, and traffic monitors characterize usage, but neither answers the fundamental question the operator actually faces---\emph{did this routing change affect the business?} Concretely, the operator must simultaneously determine the nature of the announcement (an outage, a leak, a hijack, or a benign change), whether it produced a traffic anomaly, and crucially, whether that anomaly was a direct consequence of the routing event rather than an independent phenomenon. Answering these questions requires the two signals to be joined on a common time grid and interpreted causally at minute-level latency. Unfortunately, the scientific literature has produced detectors that are highly capable on their own isolated signals~\cite{artemis}, but they emit unsorted alert streams that the operator must still arbitrate by hand. Neither the existing operational tooling nor the current scientific literature provides a unified system that answers these correlated questions together at the required minute-level latency.

The motivation for this work follows directly from this gap. We aim to build a system that bridges the routing and traffic signals on a common time grid, deterministically arbitrates the evidence from multiple detectors, and returns an auditable verdict at minute-level latency. The primary benefit of this approach is the transformation of the operator's workflow: replacing a slow, manual, and non-reproducible diagnostic process with a reproducible verdict that an operator can read in a minute rather than reconstruct over hours. Consequently, the system drastically reduces the operator's daily review load from $\mathcal{O}(10^3)$ unsorted alarms to a mere $\mathcal{O}(10)$ actionable verdicts, filtering out the overwhelming noise and surfacing only falsifiable hypotheses. The core contribution of this paper is therefore a comprehensive, end-to-end system for joint traffic-routing analysis, rather than just another isolated anomaly detector.

Realizing such a system, however, is not a matter of simply assembling existing detectors. Three technical challenges stand between the current state of the art and a practical, minute-level solution.

\paragraph{C1 --- Scale vs.\ real-time constraints.} In real operations, operators cannot afford to wait hours for an exact, post-hoc report. When an alert triggers, they need an immediate, high-level assessment to know whether it deserves urgent attention, deferring deep analysis to later. However, public BGP telemetry and traffic data are massive. Performing full-scan analysis over the complete data stream is computationally prohibitive and far too slow. The challenge is to process massive datasets and rapidly deliver an answer under strict minute-level latency.

\paragraph{C2 --- Anomaly detection on heterogeneous AS traffic.} Detecting traffic anomalies is fundamentally difficult because not all ASes are suitable for statistical detection. In real networks, the traffic of some ASes fluctuates sharply and lacks stable regularity; applying detection to them yields only noise. The challenge is to dynamically identify which ASes are genuinely detectable and to accurately judge anomalies within a minute-level window, without being overwhelmed by the inherent noise of irregular traffic.

\paragraph{C3 --- Causal attribution across routing and traffic.} The relationship between routing and traffic is not simply linear. Because routing anomalies propagate, a single fault at one AS can manifest as multiple anomalies across the routing graph, making the causal chain difficult to trace. Granger causality requires stationary residuals, which bursty BGP streams systematically violate; PC answers conditional independence, not sufficiency. Pearl's probability of necessity and sufficiency (PNS)~\cite{pearl-causality-2009} asks sufficiency directly, but BGP telemetry carries no a priori causal graph. We adopt a graph-free identification: a pre/post partition at the routing breakpoint anchors the Pearl-PNS estimator, and Granger, transfer entropy, and PC are demoted to evidence channels in a Bayesian synthesis. This demotion is the design choice that most distinguishes this system from prior work.

These three challenges directly dictate the design of BGP Tracer. We address them with three corresponding solutions.

\paragraph{S1 --- Two-stage coarse-to-fine framework (Addresses C1).} We design a ``coarse-to-fine'' two-stage strategy to balance scale and latency. The coarse screening stage rapidly filters out low-traffic or evidently normal ASes by leveraging shared caches, parallel processing, and data pruning, prioritizing the most suspicious targets. The fine-grained stage then performs detailed analysis only on these candidates, generating a comprehensive and interpretable diagnostic report. This prioritized scheduling drastically reduces computational load and guarantees the generation of minute-level assessments.

\paragraph{S2 --- Context-aware detection for regular large ASes (Addresses C2).} Since irregular ASes yield only noise, we deliberately exclude them and focus on Tier-1-level large ASes that exhibit strong periodic regularity. The system first establishes a periodic baseline for these targets, and then judges anomalies by combining current context with historical traffic patterns. This approach markedly improves detection accuracy and significantly reduces false positives caused by natural traffic volatility.

\paragraph{S3 --- Pearl-PNS-anchored causal attribution (Addresses C3).} Given the limitations of Granger causality and the PC algorithm, we adopt a graph-free Pearl-PNS identification: a pre/post partition at the routing breakpoint anchors two Pearl-PNS-anchored estimators; Granger, transfer entropy, and PC become three of five evidence channels in a Bayesian synthesis whose abstention threshold $\tau=0.34$ is the natural boundary at which the dominant hypothesis outweighs the other two combined. The core contribution is not the individual channels but the synthesis and the auditable abstention mechanism.

The main contributions of this paper are as follows:

\begin{enumerate}
\item \textbf{An agentic two-stage analysis tool.} Through the coarse-to-fine strategy, the system schedules computational resources efficiently, solving the minute-level monitoring problem under massive data scale.
\item \textbf{Context-aware anomaly detection for regular ASes.} By establishing a periodic baseline and incorporating historical context, the system effectively reduces false alarms and improves accuracy for Tier-1 ASes.
\item \textbf{A Pearl-PNS-anchored causal attribution method.} By anchoring two Pearl-PNS estimators to a graph-free pre/post partition, the system synthesises five evidence channels into an auditable Bayesian posterior with explicit abstention.
\end{enumerate}